% ---------------------------------------------------------------------
%  1  INTRODUCTION
% ---------------------------------------------------------------------
\section{Introduction}
\label{sec:intro}

The assigned task is to predict, for each of the 24 hours of the next
day, how much the price of Bitcoin will change, $\ret_h =
\mathrm{Close}_h/\mathrm{Close}_{h-1}-1$, from the hours before. We
forecast changes rather than prices because the price level is not
stationary: it rose from about \$4{,}000 (March 2020) to over \$126{,}000,
and 72\% of the test hours trade above the highest training price, so a
model of the level would have to extrapolate.

Short-horizon forecasts matter for risk management, position sizing and
the timing of trades, and Bitcoin is a demanding case: it trades 24/7,
its hourly returns are heavy-tailed, and its behaviour changes between
bull and bear markets and after events such as the launch of the US spot
ETFs in January 2024. Two classical results frame the problem. In an
efficient market, past prices should not predict the direction of future
returns~\cite{fama1970}; the size of the moves, on the other hand,
clusters in time and is predictable~\cite{engle1982,bollerslev1986}.
Financial time series also make it easy to fool oneself: overlapping
windows and random splits leak the future into the training
data~\cite{lopezdeprado2018}.

Our contributions are: (i) a leakage-free evaluation in which every
choice is made without the test block and every model is compared with a
simple rule; (ii) evidence that hourly returns are not predictable beyond
a tiny one-hour reversal, while their size is; (iii) calibrated 80\%
ranges for every hour of the next day; and (iv) a comparison of classical
models, neural networks and a zero-shot foundation model on identical
rows.
